% =========================================================
% Proof Stub: Characterization of Divergence
% Source: volume-iii/analysis/sequences/notes/sequences/notes-subsequences.tex
% Mode: standard
% =========================================================

% *** HARDER PROOF — attempt after foundational results settled ***

\clearpage
\phantomsection
\hypertarget{proof-RA-ABB-T-char-divergence}{}

\subsubsection[Characterization of Divergence]{Proof --- Characterization of Divergence}

\bigskip

\noindent
\textbf{Source.}
See \texttt{notes-subsequences.tex}.

\vspace{0.75em}

\noindent
\textbf{Goal.}
$(a_n)$ diverges iff it is unbounded, or has two subsequences with different limits, or has a subsequence diverging to $\pm\infty$.

\vspace{0.75em}

\noindent
\textbf{Logical form.}
\[
  (a_n) \text{ diverges} \Longleftrightarrow (a_n) \text{ unbounded} \vee \exists (a_{n_k}) \to L,\; (a_{m_j}) \to M,\; L \ne M \vee \ldots
\]

\vspace{0.75em}

\noindent
\textbf{Proof strategy.}
Prove by contrapositive/cases: if convergent then bounded and all subsequences share the same limit. For the converse cases, exhibit a specific divergent witness.

\vspace{0.75em}

\noindent
\textbf{Proof.}
\begin{proof}
\Claim $(a_n)$ diverges iff it is unbounded, or has two subsequences with different limits, or has a subsequence diverging to $\pm\infty$.

\Given 

\Goal 

\Strategy Prove by contrapositive/cases: if convergent then bounded and all subsequences share the same limit. For the converse cases, exhibit a specific divergent witness.

\medskip

% --- (⇐) each case produces divergence directly or by contradiction with uniqueness of limits ---
% --- (⇒) if convergent → bounded (conv-bounded thm) and all subsequences → same limit ---

\end{proof}

\begin{remark*}[Proof shape]
Combine: convergent $\Rightarrow$ bounded $+$ unique limit, and check all three failure modes go the other way.
\end{remark*}

\begin{remark*}[Dependencies]
\begin{itemize}
  \item Convergent sequences are bounded.
  \item Subsequences inherit limits.
  \item Uniqueness of limits.
\end{itemize}
\end{remark*}
